% ===========================================================================
%  RAPPORT TECHNIQUE — PROJET 3
%  Conception et mise en oeuvre d'une pipeline DevSecOps augmentee par IA
%  pour les applications de l'ecosysteme JavaScript (Express.js)
%
%  Module  : Securite des applications et des services
%  Filiere : Cycle ingenieur IRISI 3
%  Etablis.: Faculte des Sciences et Techniques Marrakech (UCA)
%  Annee   : 2026/2027
%
%  COMPILATION : pdfLaTeX, deux passes minimum
%                (voir README.md pour les instructions Overleaf)
% ===========================================================================

\documentclass[11pt,a4paper,oneside]{report}

% ===========================================================================
%  PREAMBULE
%  Paquets volontairement limites a des paquets standards disponibles
%  sur toute distribution TeX Live / MiKTeX recente et sur Overleaf.
% ===========================================================================

% ---- Encodage et langue ---------------------------------------------------
\usepackage[utf8]{inputenc}
\usepackage[T1]{fontenc}
\usepackage[french]{babel}
\usepackage{lmodern}

% ---- Mise en page ---------------------------------------------------------
\usepackage[a4paper,margin=2.5cm,headheight=14pt]{geometry}
\usepackage{fancyhdr}
\usepackage{setspace}
\onehalfspacing

% ---- Tableaux et figures --------------------------------------------------
\usepackage{graphicx}
\usepackage{booktabs}
\usepackage{longtable}
\usepackage{array}
\usepackage{tabularx}
\usepackage{multirow}
\usepackage{float}
\usepackage{caption}

% ---- Couleurs et encadres -------------------------------------------------
\usepackage{xcolor}
\usepackage{framed}

% ---- Listes ---------------------------------------------------------------
\usepackage{enumitem}

% ---- Code source ----------------------------------------------------------
\usepackage{listings}

% ---- Symboles -------------------------------------------------------------
\usepackage{amssymb}
\usepackage{pifont}

% ---- Liens hypertexte (a charger en dernier) ------------------------------
\usepackage[hidelinks,breaklinks=true]{hyperref}

% ===========================================================================
%  COULEURS DU DOCUMENT
% ===========================================================================
\definecolor{bleufonce}{RGB}{11,61,92}
\definecolor{bleuclair}{RGB}{244,247,249}
\definecolor{grisbord}{RGB}{180,180,180}
\definecolor{orangeattention}{RGB}{217,145,0}
\definecolor{jauneclair}{RGB}{255,248,230}
\definecolor{vertfonce}{RGB}{45,107,45}
\definecolor{vertclair}{RGB}{234,244,234}
\definecolor{rougeplaceholder}{RGB}{176,0,32}
\definecolor{griscode}{RGB}{245,245,245}
\definecolor{commentairevert}{RGB}{0,110,0}
\definecolor{chainemauve}{RGB}{140,0,140}

% ===========================================================================
%  EN-TETES ET PIEDS DE PAGE
% ===========================================================================
\pagestyle{fancy}
\fancyhf{}
\fancyhead[L]{\small\nouppercase{\leftmark}}
\fancyhead[R]{\small\thepage}
\renewcommand{\headrulewidth}{0.4pt}

\fancypagestyle{plain}{%
  \fancyhf{}
  \fancyhead[R]{\small\thepage}
  \renewcommand{\headrulewidth}{0pt}
}

% ===========================================================================
%  STYLE DU CODE SOURCE
% ===========================================================================
\lstdefinestyle{stylegeneral}{
  backgroundcolor=\color{griscode},
  basicstyle=\ttfamily\scriptsize,
  keywordstyle=\color{bleufonce}\bfseries,
  commentstyle=\color{commentairevert}\itshape,
  stringstyle=\color{chainemauve},
  numberstyle=\tiny\color{gray},
  numbers=left,
  numbersep=8pt,
  frame=single,
  rulecolor=\color{grisbord},
  breaklines=true,
  breakatwhitespace=false,
  showstringspaces=false,
  tabsize=2,
  captionpos=b,
  xleftmargin=14pt,
  framexleftmargin=14pt,
  literate=%
    {é}{{\'e}}1 {è}{{\`e}}1 {ê}{{\^e}}1 {ë}{{\"e}}1
    {à}{{\`a}}1 {â}{{\^a}}1 {ô}{{\^o}}1 {î}{{\^i}}1
    {ï}{{\"i}}1 {ù}{{\`u}}1 {û}{{\^u}}1 {ç}{{\c c}}1
    {É}{{\'E}}1 {È}{{\`E}}1 {À}{{\`A}}1
}
\lstset{style=stylegeneral}

\lstdefinelanguage{yamlperso}{
  keywords={name,on,jobs,steps,uses,with,run,runs-on,permissions,
            concurrency,if,needs,env,repos,hooks,rev,id,args},
  sensitive=true,
  comment=[l]{\#},
  morestring=[b]",
  morestring=[b]'
}

% ===========================================================================
%  COMMANDES PERSONNALISEES
% ===========================================================================

% --- Reference a une pratique DevSecOps : \pratique{P1}{SAST}
\newcommand{\pratique}[2]{\textbf{#1}~---~#2}

% --- Identifiant de pratique seul : \pid{P1}
\newcommand{\pid}[1]{\textbf{#1}}

% --- Exigence ASVS : \asvs{8.2.2}
\newcommand{\asvs}[1]{\mbox{\textsc{asvs}~v5.0.0--#1}}

% --- Exigence ASVS courte (dans un tableau) : \asvsc{8.2.2}
\newcommand{\asvsc}[1]{\mbox{\textsc{asvs}~#1}}

% --- Risque OWASP LLM : \llm{01}
\newcommand{\llm}[1]{\mbox{\textsc{llm}#1}}

% --- Marqueurs d'etat
\newcommand{\fait}{\textcolor{vertfonce}{\ding{51}}}
\newcommand{\nonfait}{\textcolor{grisbord}{\ding{55}}}
\newcommand{\encours}{\textcolor{orangeattention}{$\circlearrowright$}}
\newcommand{\bloque}{\textcolor{rougeplaceholder}{\ding{47}}}

% ===========================================================================
%  EMPLACEMENT DE CAPTURE D'ECRAN MANQUANTE
%  Affiche un encadre bien visible dans le PDF pour ne rien oublier.
%  Usage : \captureaplacer{figures/nom.png}{description de la capture}
% ===========================================================================
\newcommand{\captureaplacer}[2]{%
  \begin{figure}[H]
  \centering
  \setlength{\fboxrule}{1.2pt}%
  \setlength{\fboxsep}{10pt}%
  \textcolor{rougeplaceholder}{%
    \fbox{%
      \begin{minipage}{0.86\textwidth}
        \centering
        \color{rougeplaceholder}%
        \textbf{\large CAPTURE D'ÉCRAN À AJOUTER}\\[6pt]
        \color{black}%
        \textbf{Fichier attendu :} \texttt{#1}\\[4pt]
        \textbf{Contenu :} #2\\[8pt]
        {\small\itshape Déposer l'image dans \texttt{figures/} sous ce nom
        exact, puis remplacer \texttt{\textbackslash captureaplacer}
        par \texttt{\textbackslash captureprete} dans le fichier source.}
        \vspace{4pt}
      \end{minipage}%
    }%
  }
  \end{figure}
}

% ===========================================================================
%  CAPTURE FOURNIE
%  Usage : \captureprete{figures/nom.png}{Légende de la figure}
%
%  ROBUSTE PAR CONCEPTION : si le fichier image est absent, la commande
%  affiche l'encadre de rappel au lieu de provoquer une erreur de
%  compilation. Un fichier oublie ne peut donc pas casser le document.
% ===========================================================================
\newcommand{\captureprete}[2]{%
  \IfFileExists{#1}{%
    \begin{figure}[H]
      \centering
      \includegraphics[width=0.92\textwidth]{#1}
      \caption{#2}
    \end{figure}
  }{%
    \captureaplacer{#1}{#2}%
  }%
}

% Variante en largeur reduite, meme robustesse
\newcommand{\capturepretedemi}[2]{%
  \IfFileExists{#1}{%
    \begin{figure}[H]
      \centering
      \includegraphics[width=0.60\textwidth]{#1}
      \caption{#2}
    \end{figure}
  }{%
    \captureaplacer{#1}{#2}%
  }%
}

% ===========================================================================
%  ENCADRES THEMATIQUES
% ===========================================================================

% --- Encadre « decision d'architecture »
\newenvironment{decision}[1]{%
  \begin{framed}
  \noindent\textcolor{bleufonce}{\textbf{Décision --- #1}}\par\vspace{2pt}
  \small
}{%
  \end{framed}
}

% --- Encadre « constat / point d'attention »
\newenvironment{constat}[1]{%
  \begin{framed}
  \noindent\textcolor{orangeattention}{\textbf{Constat --- #1}}\par\vspace{2pt}
  \small
}{%
  \end{framed}
}

% --- Encadre « verification effectuee »
\newenvironment{verification}[1]{%
  \begin{framed}
  \noindent\textcolor{vertfonce}{\textbf{Vérification --- #1}}\par\vspace{2pt}
  \small
}{%
  \end{framed}
}

% ===========================================================================
%  CHAINE DE DEMONSTRATION EXIGEE PAR LE SUJET
%  Pratique -> Outil -> Integration -> Detection -> Resultat -> Correction
% ===========================================================================
\newcommand{\chainedemo}[6]{%
  \begin{center}
  \small
  \begin{tabular}{@{}>{\bfseries}p{3.1cm}p{10.4cm}@{}}
    \toprule
    Pratique      & #1 \\
    Outil         & #2 \\
    Intégration   & #3 \\
    Détection     & #4 \\
    Résultat      & #5 \\
    Correction    & #6 \\
    \bottomrule
  \end{tabular}
  \end{center}
}

% ===========================================================================
%  DIVERS
% ===========================================================================
\setlist[itemize]{itemsep=2pt,topsep=4pt,parsep=0pt}
\setlist[enumerate]{itemsep=2pt,topsep=4pt,parsep=0pt}

\captionsetup{font=small,labelfont=bf}

% Eviter les coupures de mots trop agressives
\tolerance=1000
\emergencystretch=2em

% Profondeur de la table des matieres
\setcounter{tocdepth}{2}
\setcounter{secnumdepth}{3}


% ---------------------------------------------------------------------------
\begin{document}
% ---------------------------------------------------------------------------

% ===========================================================================
%  PAGE DE GARDE
% ===========================================================================
\begin{titlepage}
\centering

{\large\scshape Université Cadi Ayyad}\\[2pt]
{\large\scshape Faculté des Sciences et Techniques de Marrakech}\\[2pt]
{\large\scshape Département d'Informatique}\\[4pt]
\rule{\textwidth}{1pt}\\[1.2cm]

{\scshape\large Cycle Ingénieur --- IRISI 3}\\[4pt]
{\scshape Module : Sécurité des applications et des services}\\[1.6cm]

{\Large\bfseries Projet du module --- Projet 3}\\[1.0cm]

\rule{0.85\textwidth}{0.6pt}\\[0.5cm]
{\huge\bfseries Conception et mise en œuvre d'une pipeline DevSecOps
augmentée par IA}\\[0.5cm]
{\LARGE pour les applications de l'écosystème JavaScript}\\[0.3cm]
{\large\itshape Cas d'étude : place de marché d'articles vintage}\\[0.4cm]
\rule{0.85\textwidth}{0.6pt}\\[1.8cm]

\begin{minipage}{0.45\textwidth}
\begin{flushleft}
\textbf{Réalisé par :}\\[4pt]
ELKHANFAF Mohamed\\
ELFADIL Hanane\\
ERRAGUIBI Abdelilah\\
HALLOUMI Anas
\end{flushleft}
\end{minipage}
\hfill
\begin{minipage}{0.45\textwidth}
\begin{flushright}
\textbf{Encadré par :}\\[4pt]
% TODO : nom de l'enseignant encadrant
\textit{[Nom de l'encadrant]}
\end{flushright}
\end{minipage}

\vfill

{\small\itshape Document rédigé en parallèle de la réalisation du projet.}\\[2pt]
{\small Année universitaire 2026/2027}

\end{titlepage}

% ===========================================================================
%  RESUME
% ===========================================================================
\chapter*{Résumé}
\addcontentsline{toc}{chapter}{Résumé}
\markboth{Résumé}{Résumé}

Ce rapport présente la conception et la mise en œuvre d'une pipeline DevSecOps
augmentée par intelligence artificielle, appliquée à une application web de
l'écosystème JavaScript. Le cas d'étude retenu est une place de marché
d'articles vintage, délibérément réduite en nombre de fonctionnalités mais
choisie pour la richesse de sa surface d'attaque : authentification,
autorisation fondée sur la propriété des ressources, téléversement de fichiers,
logique métier à valeur marchande et intégration d'un modèle de langage.

La démarche repose sur deux référentiels. La taxonomie des pratiques DevSecOps
de Prates et Pereira~\cite{prates2025} fournit la liste des pratiques à mettre
en œuvre, dont les onze pratiques de niveau Tier~1 constituent le socle
obligatoire du projet. Le référentiel OWASP ASVS~5.0~\cite{owaspasvs5} fournit
les exigences de sécurité applicative vérifiables et le niveau d'assurance visé.

L'originalité de la démarche adoptée tient à deux choix structurants. D'une
part, la gouvernance de sécurité est mise en place \emph{avant} toute écriture
de code applicatif, de sorte que chaque contribution ultérieure soit contrôlée
dès le premier envoi. D'autre part, la démonstration de l'efficacité de la
pipeline repose sur l'introduction volontaire et contrôlée de vulnérabilités :
une application exempte de défaut ne produit que des rapports vides, lesquels
ne démontrent la détection d'aucune pratique.

\vspace{0.6cm}
\noindent\textbf{Mots-clés :} DevSecOps, sécurité applicative, intégration
continue, OWASP ASVS, SAST, DAST, SCA, modélisation des menaces,
Express.js, intelligence artificielle générative.

% ===========================================================================
%  TABLES
% ===========================================================================
\tableofcontents
\listoffigures
\listoftables

% ===========================================================================
%  CHAPITRES
% ===========================================================================

% ===========================================================================
\chapter{Introduction}
\label{chap:introduction}
% ===========================================================================

\section{Contexte}

Le rythme de livraison des logiciels s'est considérablement accéléré. Les
organisations informatiques sont soumises à des cycles de publication toujours
plus courts, tout en devant satisfaire des exigences non fonctionnelles
croissantes, en particulier en matière de sécurité et de conformité. L'adoption
de DevOps a constitué la réponse dominante à cette pression, en automatisant la
chaîne de construction et de déploiement~\cite{prates2025}.

Cette accélération a toutefois créé un déséquilibre. Lorsque la sécurité reste
une activité de fin de cycle, elle devient soit un goulot d'étranglement, soit
une étape escamotée. DevSecOps désigne l'intégration des pratiques de sécurité
dans DevOps~\cite{myrbakken2017}, selon le principe du \emph{shift security to
left} : faire intervenir la sécurité dès les premières étapes du cycle de
développement plutôt qu'au moment de la mise en production.

Le présent projet s'inscrit dans ce cadre. Il ne s'agit pas de développer une
application riche en fonctionnalités, puis de l'auditer. Il s'agit de concevoir
une chaîne de production logicielle dans laquelle les contrôles de sécurité sont
automatisés, outillés, et dont l'efficacité est démontrable.

\section{Problématique}

Le sujet du module formule une exigence précise quant à la nature de la
démonstration attendue : il ne suffit pas de présenter une liste d'outils, mais
d'établir pour chaque mécanisme de sécurité la chaîne complète

\begin{center}
\textbf{Pratique de sécurité $\rightarrow$ Outil $\rightarrow$ Intégration dans
la pipeline $\rightarrow$ Détection $\rightarrow$ Résultat $\rightarrow$ Action
corrective}
\end{center}

Cette exigence soulève une difficulté méthodologique qui a structuré l'ensemble
du projet et qu'il convient d'énoncer d'emblée.

\begin{constat}{une application sécurisée ne démontre rien}
Si le code produit est exempt de défaut dès le premier commit, chaque analyseur
de la pipeline retourne un rapport vide. Or un rapport vide ne démontre les
étapes \emph{Détection}, \emph{Résultat} et \emph{Action corrective} pour
aucune pratique.

La capacité à démontrer l'efficacité d'un contrôle suppose donc de disposer de
détections réelles. Nous avons choisi de les produire de manière volontaire et
contrôlée, selon la démarche décrite au chapitre~\ref{chap:vulnerabilites}.
\end{constat}

Une seconde difficulté concerne la nature même des pratiques à mettre en œuvre.
Comme le montre le chapitre~\ref{chap:fondements}, plusieurs des pratiques de
niveau Tier~1 ne sont pas outillables : elles relèvent du processus et de
l'organisation. Leur mise en œuvre ne produit pas de rapport d'outil mais des
artefacts documentaires et des configurations de plate-forme. Les traiter comme
des étapes de pipeline aurait conduit à un inventaire incomplet.

\section{Objectifs}

Le projet poursuit les objectifs suivants.

\begin{enumerate}
  \item \textbf{Identifier et mettre en œuvre les onze pratiques DevSecOps de
  niveau Tier~1} telles que définies par Prates et Pereira~\cite{prates2025},
  en automatisant leur exécution autant que possible.

  \item \textbf{Mettre en œuvre un sous-ensemble justifié de pratiques de
  niveau Tier~2}, en précisant pour chacune son objectif, l'outil retenu, son
  intégration dans la pipeline, les résultats obtenus et sa contribution à la
  sécurité.

  \item \textbf{Développer une application} comportant un front-end web, une
  application mobile, un ensemble d'API REST et un back-end métier, dont la
  sécurité est intégrée au cycle de développement.

  \item \textbf{Prendre en compte les catégories de cyber-risques} issues des
  référentiels OWASP : Top~10:2025 pour le web~\cite{owasptop10}, API
  Security Top~10~\cite{owaspapi}, Mobile Top~10~\cite{owaspmobile} et
  Top~10 for LLM Applications~\cite{owaspllm}.

  \item \textbf{Augmenter la pipeline par l'intelligence artificielle}, tant
  dans l'application que dans la chaîne d'intégration continue elle-même.

  \item \textbf{Démontrer l'efficacité de chaque contrôle} selon la chaîne de
  démonstration exigée, preuves à l'appui.

  \item \textbf{Analyser l'impact des contrôles de sécurité sur la performance
  de la pipeline}, et documenter les compromis retenus.
\end{enumerate}

\section{Périmètre et stratégie de réalisation}

\subsection{Priorité accordée à la profondeur plutôt qu'à l'étendue}

Le sujet précise que la priorité d'évaluation porte sur la qualité de la
démarche de cybersécurité, l'automatisation des contrôles et la capacité à
démontrer l'efficacité de la pipeline, plutôt que sur le nombre de
fonctionnalités développées.

Cette indication a été prise au sérieux. Le périmètre fonctionnel de
l'application est volontairement restreint, et chaque fonctionnalité retenue
l'a été parce qu'elle crée une surface de sécurité que nous devons ensuite
défendre, outiller et démontrer. Les fonctionnalités n'apportant aucune
catégorie de risque nouvelle ont été écartées, notamment le traitement réel des
paiements, l'envoi de courriels et l'authentification fédérée.

\subsection{Réalisation séquentielle du web puis du mobile}

L'application mobile, exigée par le sujet, a été délibérément reportée à une
seconde phase, postérieure à l'achèvement complet de l'application web et de sa
pipeline.

\begin{decision}{report de l'application mobile}
Une application web complète dotée d'une pipeline entièrement outillée et
documentée a davantage de valeur, au regard des critères d'évaluation énoncés,
que deux clients partiellement achevés partageant une pipeline superficielle.

Ce report n'est pas un abandon : l'architecture a été conçue dès l'origine
autour d'une API indifférente au client, de sorte que l'ajout du client mobile
ne nécessite aucune reprise côté serveur. Les travaux restants sont recensés
et planifiés.
\end{decision}

\subsection{Rédaction du rapport en parallèle de la réalisation}

Le présent rapport est rédigé au fur et à mesure de l'avancement, chaque phase
achevée étant documentée immédiatement. Ce choix répond à une contrainte
pratique : le raisonnement qui conduit à une décision technique est la première
chose qui s'efface. Les registres de décisions d'architecture
(\emph{Architecture Decision Records}) tenus tout au long du projet constituent
la matière première des sections de justification de ce document.

Cette démarche est par ailleurs cohérente avec les exigences de documentation
introduites par l'ASVS~5.0~\cite{owaspasvs5}, qui demande que certaines
décisions de sécurité soient explicitement consignées pour pouvoir être
vérifiées.

\section{Organisation du rapport}

Le chapitre~\ref{chap:fondements} présente les fondements de la démarche : le
passage de DevOps à DevSecOps, la taxonomie des pratiques retenue et les
référentiels de sécurité appliqués. Il identifie en particulier les onze
pratiques de niveau Tier~1.

Les chapitres~\ref{chap:archi-app} et~\ref{chap:archi-pipeline} décrivent
respectivement l'architecture de l'application et celle de la pipeline
DevSecOps.

Le chapitre~\ref{chap:phase0} documente la première phase de réalisation : la
mise en place de la gouvernance de sécurité du dépôt, antérieure à toute
écriture de code applicatif.

Les chapitres suivants documentent les phases ultérieures au fur et à mesure de
leur achèvement : modélisation des menaces, développement sécurisé, tests de
sécurité automatisés, étages de sécurité de la pipeline, conteneurisation,
supervision, augmentation par l'IA.

Le chapitre~\ref{chap:vulnerabilites} présente les démonstrations de
vulnérabilités volontaires et les corrections apportées. Le
chapitre~\ref{chap:performance} analyse l'impact des contrôles sur la
performance de la pipeline. Le chapitre~\ref{chap:limites} expose les limites du
travail et les perspectives.

% ===========================================================================
\chapter{Fondements et référentiels}
\label{chap:fondements}
% ===========================================================================

\section{De DevOps à DevSecOps}

DevOps est défini comme une approche organisationnelle qui met l'accent sur la
collaboration transversale au sein et entre les équipes de développement
logiciel, afin d'exploiter des systèmes résilients et d'accélérer la livraison
des changements~\cite{prates2025}. Ses principes sont fréquemment résumés par
l'acronyme CAMS : culture, automatisation, mesure et partage.

DevSecOps étend ces principes en y intégrant la sécurité. Prates et Pereira
montrent que chacun des quatre principes CAMS admet une déclinaison sécurité :
la culture s'étend à la collaboration avec l'équipe sécurité, l'automatisation
englobe l'automatisation des contrôles de sécurité, la mesure intègre des
métriques de menaces et de vulnérabilités, et le partage inclut la mise en
commun des problématiques de sécurité~\cite{prates2025}.

Les auteurs considèrent DevSecOps comme une extension naturelle de DevOps,
position également retenue par la norme IEEE~2675~\cite{ieee2675}.

Le principe structurant est celui du \emph{shift security to left} : inclure les
pratiques de sécurité dès les étapes les plus précoces du cycle de
développement. Ce principe a une conséquence directe sur l'ordonnancement des
travaux, exposée au chapitre~\ref{chap:phase0}.

\section{Le cycle de vie DevSecOps}

Prates et Pereira identifient les phases du flux DevSecOps et y rattachent
chaque pratique~\cite{prates2025}. Les phases sont les suivantes :
\emph{Create}, \emph{Plan}, \emph{Adapt}, \emph{Verify},
\emph{Preproduction}, \emph{Release}, \emph{Predict}, \emph{Respond},
\emph{Detect} et \emph{Configure}.

Cette cartographie présente un intérêt pratique immédiat : elle indique
\emph{où} dans la chaîne chaque pratique doit intervenir. Elle a servi de
squelette à l'architecture de notre pipeline, présentée au
chapitre~\ref{chap:archi-pipeline}.

\section{La taxonomie des pratiques de Prates et Pereira}

\subsection{Méthode et portée de l'étude}

L'étude de référence du module est une revue de littérature multivocale
(\emph{Multivocal Literature Review}) portant sur 103 publications
primaires~\cite{prates2025}. Contrairement à une revue systématique classique
limitée à la littérature académique, une revue multivocale intègre également la
littérature grise : rapports techniques d'entreprises, pages web, livres
blancs.

La répartition des sources est quasiment équilibrée, avec 49\% de littérature
blanche et 51\% de littérature grise, ce que les auteurs interprètent comme une
convergence entre la recherche académique et la pratique industrielle.

L'étude identifie \textbf{39 pratiques DevSecOps distinctes}.

\subsection{Le classement par niveaux}

Les auteurs classent ces 39 pratiques en trois niveaux selon le nombre de
publications les mentionnant~\cite{prates2025} :

\begin{itemize}
  \item \textbf{Tier 1 --- Core} : pratiques mentionnées dans \textbf{30
  publications ou plus}. Elles constituent les pratiques fondamentales de
  DevSecOps.
  \item \textbf{Tier 2 --- Important} : pratiques mentionnées dans
  \textbf{10 à 29} publications. Également considérées comme importantes.
  \item \textbf{Tier 3 --- Nice to have} : pratiques mentionnées dans
  \textbf{moins de 10} publications. Interprétées comme souhaitables mais non
  cruciales pour une mise en œuvre de DevSecOps.
\end{itemize}

Il importe de souligner que le niveau d'une pratique reflète son degré
d'établissement dans la littérature, et \textbf{non} son caractère automatisable
ni sa position dans la pipeline. Cette distinction a des conséquences directes,
développées à la section~\ref{sec:pratiques-non-outillables}.

\subsection{Les onze pratiques de niveau Tier 1}

Le tableau~\ref{tab:tier1} présente les onze pratiques de niveau Tier~1, qui
constituent le socle obligatoire du projet.

\begin{table}[H]
\centering
\caption{Les onze pratiques DevSecOps de niveau Tier~1 (d'après
Prates et Pereira~\cite{prates2025}, tableau~7)}
\label{tab:tier1}
\small
\begin{tabular}{@{}llcl@{}}
\toprule
\textbf{Id.} & \textbf{Pratique} & \textbf{Mentions} & \textbf{Phase du cycle} \\
\midrule
P1  & SAST                        & 65 & Verify \\
P2  & DAST                        & 58 & Verify \\
P3  & Automated security testing  & 56 & Verify \\
P4  & Vulnerability scanning      & 44 & Verify / Preproduction \\
P5  & Continuous monitoring       & 41 & Predict / Detect \\
P6  & Continuous integration      & 41 & Adapt \\
P7  & Threat modelling            & 40 & Plan / Predict \\
P8  & Shift left security         & 36 & Create \\
P9  & Collaboration               & 36 & Create \\
P10 & Continuous deployment       & 34 & Release \\
P11 & Code review                 & 33 & Verify \\
\bottomrule
\end{tabular}
\end{table}

Les définitions retenues, reformulées d'après l'article source, sont les
suivantes.

\begin{description}[leftmargin=1.2cm,style=nextline]
  \item[\pid{P1} SAST] Analyse du code source, du bytecode ou du binaire d'une
  application afin d'identifier des vulnérabilités, sans exécuter le programme.

  \item[\pid{P2} DAST] Test d'une application en cours d'exécution, vue de
  l'extérieur, par envoi de requêtes et analyse des réponses, en simulant des
  attaques réelles.

  \item[\pid{P3} Automated security testing] Recours à des outils et processus
  automatisés pour vérifier systématiquement la sécurité à différentes étapes
  du développement. Couvre les tests d'API, les tests d'applications web et les
  tests unitaires orientés sécurité. Cette pratique est \emph{distincte} de
  SAST et de DAST.

  \item[\pid{P4} Vulnerability scanning] Identification automatisée et continue
  des vulnérabilités du logiciel \emph{et} de l'infrastructure, par
  confrontation à des bases de CVE connues. L'article précise que cette pratique
  englobe l'analyse des dépendances logicielles et celle des images de
  conteneurs.

  \item[\pid{P5} Continuous monitoring] Collecte et analyse automatisées et
  continues des données de performance et de sécurité des systèmes en
  fonctionnement, procurant une visibilité en temps réel.

  \item[\pid{P6} Continuous integration] Intégration automatique et continue
  des modifications de code dans un dépôt partagé, avec construction, test,
  validation et contrôles de sécurité automatisés.

  \item[\pid{P7} Threat modelling] Processus systématique et proactif
  d'identification et d'évaluation des menaces, par analyse de la conception, de
  l'architecture et des fonctionnalités du système.

  \item[\pid{P8} Shift left security] Intégration des considérations de sécurité
  le plus tôt possible dans le cycle de développement, plutôt qu'à la fin.

  \item[\pid{P9} Collaboration] Collaboration transversale entre les équipes de
  développement, d'exploitation et de sécurité, par décloisonnement, afin de
  créer une culture consciente de la sécurité.

  \item[\pid{P10} Continuous deployment] Construction, test et déploiement
  automatiques et continus des modifications vers la production, les contrôles
  de sécurité étant partie intégrante de la chaîne automatisée.

  \item[\pid{P11} Code review] Examen systématique du code par des pairs ou des
  experts afin d'identifier défauts, vulnérabilités et problèmes de qualité
  avant fusion.
\end{description}

\subsection{Pratiques de niveau Tier 2}

Le tableau~\ref{tab:tier2} recense les vingt pratiques de niveau Tier~2 et
indique celles retenues pour ce projet.

\begin{table}[H]
\centering
\caption{Pratiques de niveau Tier~2 et sélection retenue}
\label{tab:tier2}
\small
\begin{tabular}{@{}llcl@{}}
\toprule
\textbf{Id.} & \textbf{Pratique} & \textbf{Mentions} & \textbf{Retenue} \\
\midrule
P12 & Infrastructure as code         & 28 & \fait{} oui \\
P13 & SCA                            & 27 & \fait{} oui \\
P14 & Security training              & 27 & \fait{} oui \\
P15 & VCS                            & 26 & \fait{} oui \\
P16 & Penetration testing            & 24 & \fait{} oui \\
P17 & Continuous compliance          & 22 & \fait{} oui \\
P18 & Logging                        & 21 & \fait{} oui \\
P19 & IAST                           & 18 & \nonfait{} non \\
P20 & Secure coding                  & 18 & \fait{} oui \\
P21 & Configuration management       & 17 & optionnel \\
P22 & Environment hardening          & 16 & \fait{} oui \\
P23 & Risk assessment                & 15 & \fait{} oui \\
P24 & Security requirement gathering & 13 & \fait{} oui \\
P25 & RASP                           & 13 & \nonfait{} non \\
P26 & Security policies              & 12 & optionnel \\
P27 & Patch management               & 12 & \fait{} oui \\
P28 & Security incident management   & 11 & \fait{} oui \\
P29 & IAM                            & 11 & optionnel \\
P30 & Secrets management             & 11 & \fait{} oui \\
P31 & Security auditing              & 11 & \fait{} oui \\
\bottomrule
\end{tabular}
\end{table}

\begin{constat}{trois pratiques fréquemment supposées Tier 1 n'en sont pas}
L'analyse de la source corrige trois idées reçues qu'il importe de ne pas
reproduire dans ce rapport :

\begin{itemize}
  \item \textbf{SCA} (analyse de composition logicielle) est une pratique de
  niveau Tier~2 (\pid{P13}, 27 mentions), et non Tier~1. La pratique de niveau
  Tier~1 qui couvre l'analyse des dépendances et des images de conteneurs est
  \pid{P4} \emph{Vulnerability scanning}.
  \item \textbf{Secrets management} est également de niveau Tier~2
  (\pid{P30}, 11 mentions).
  \item \textbf{Infrastructure as code} est de niveau Tier~2 (\pid{P12},
  28 mentions), à deux mentions du seuil de 30.
\end{itemize}

Cette distinction a une conséquence favorable sur le projet. Le sujet exige par
ailleurs l'emploi d'outils SAST, SCA, \emph{secret scanning} et DAST. SAST et
DAST relevant du Tier~1 et SCA et \emph{secrets management} du Tier~2, leur
mise en œuvre satisfait simultanément une exigence obligatoire du sujet et
ouvre droit à la valorisation des pratiques Tier~2.
\end{constat}

\subsection{Pratiques de niveau Tier 3}

Pour mémoire, les huit pratiques de niveau Tier~3 sont : \emph{Security
champions} (8 mentions), \emph{Red team} (8), \emph{WAF} (8),
\emph{Alerting} (7), \emph{Continuous reporting} (6), \emph{Software
licensing} (3), \emph{Chaos engineering} (3) et \emph{Bug bounties} (2).

Trois d'entre elles sont obtenues à coût marginal quasi nul dans notre
architecture et ont été retenues à ce titre : \emph{Alerting}, fourni par la
pile de supervision ; \emph{WAF}, par un reverse-proxy en frontal ; et
\emph{Software licensing}, par une option de l'outil d'analyse de composition.

\section{Pratiques outillables et pratiques de processus}
\label{sec:pratiques-non-outillables}

Une lecture attentive du tableau~\ref{tab:tier1} fait apparaître un point
méthodologique essentiel : \textbf{quatre des onze pratiques de niveau Tier~1 ne
sont pas outillables} au sens d'un analyseur produisant un rapport.

\begin{table}[H]
\centering
\caption{Nature des pratiques Tier~1 et nature des preuves correspondantes}
\label{tab:nature-preuves}
\small
\begin{tabular}{@{}p{3.4cm}p{4.6cm}p{5.6cm}@{}}
\toprule
\textbf{Nature} & \textbf{Pratiques} & \textbf{Nature de la preuve} \\
\midrule
Automatisable en
pipeline
& \pid{P1} SAST, \pid{P2} DAST,
\pid{P3} tests automatisés,
\pid{P4} scan de vulnérabilités
& Rapport d'outil, journal
d'exécution, échec de build \\
\addlinespace
Automatisable
hors analyse
& \pid{P6} intégration continue,
\pid{P10} déploiement continu
& Définition de la pipeline,
historique d'exécutions \\
\addlinespace
Processus et
organisation
& \pid{P7} modélisation des menaces,
\pid{P8} shift left,
\pid{P9} collaboration,
\pid{P11} revue de code
& Document, configuration de
plate-forme, fil de discussion
de \emph{pull request} \\
\addlinespace
Exploitation
& \pid{P5} supervision continue
& Tableau de bord, alerte
déclenchée \\
\bottomrule
\end{tabular}
\end{table}

Cette répartition explique la formulation du sujet, qui demande d'automatiser
l'exécution des pratiques \emph{autant que possible}. Pour les pratiques de
processus, la preuve attendue n'est pas un rapport d'analyse mais l'artefact
lui-même : le modèle de menaces, les règles de protection de branche, le fichier
d'attribution de propriété du code, ou un fil de revue montrant une
vulnérabilité discutée puis corrigée.

\section{Référentiels de sécurité appliqués}

\subsection{OWASP ASVS 5.0 comme référentiel de vérification}

La taxonomie de Prates et Pereira indique \emph{quelles pratiques} mettre en
œuvre dans la chaîne de production. Elle ne dit rien des propriétés de sécurité
que l'application doit elle-même présenter. Ces deux questions sont distinctes et
appellent deux référentiels distincts.

L'ASVS~5.0~\cite{owaspasvs5} fournit le second. Il définit environ 350 exigences
de sécurité applicative réparties en 17 chapitres, organisées en trois niveaux
de vérification. L'ASVS précise explicitement qu'il ne prescrit ni les activités
du cycle de développement, ni la manière dont l'application doit être construite
par une pipeline CI/CD : il spécifie les résultats de sécurité à atteindre au
sein du produit. Les deux référentiels sont donc complémentaires et non
redondants.

Deux apports de l'ASVS~5.0 ont particulièrement orienté le projet.

Le premier est la notion de \textbf{décision de sécurité documentée}. L'ASVS
introduit des exigences de documentation au motif que certains contrôles
dépendent trop du contexte applicatif pour être prescrits de manière générique,
et que sans décision écrite les exigences de mise en œuvre correspondantes ne
sont tout simplement pas vérifiables.

Le second est la position de l'ASVS sur les limites de l'outillage automatisé.
Le référentiel indique que les outils SAST et DAST, correctement intégrés,
peuvent détecter des problèmes de sécurité, mais qu'ils ne sont pas en mesure de
vérifier un grand nombre des exigences les plus complexes, en particulier celles
relatives à la logique métier et au contrôle d'accès. Il en conclut que des
vérifications spécifiques à l'application doivent être écrites, de nature
comparable aux tests unitaires et d'intégration.

\begin{constat}{convergence des deux référentiels}
Cette position de l'ASVS rejoint exactement la pratique \pid{P3}
\emph{Automated security testing}, troisième pratique la plus citée de l'étude
de Prates et Pereira avec 56 mentions, et définie comme distincte de SAST et de
DAST.

Les deux sources convergent donc vers un même livrable : une suite de tests de
sécurité spécifiques à l'application, portant notamment sur l'autorisation et
la logique métier. Ce livrable fait l'objet du chapitre~\ref{chap:tests}.
\end{constat}

\subsection{Niveau de vérification ASVS retenu}

L'ASVS~5.0 définit trois niveaux. Le niveau~1 représente environ 20\% des
exigences et constitue la première couche de défense. Le niveau~2 correspond au
niveau que la majorité des applications devraient viser. Le niveau~3 s'adresse
aux applications exigeant la plus haute assurance.

Le référentiel encourage explicitement les organisations à produire une version
adaptée à leur domaine, en écartant les chapitres non pertinents, à condition
que le périmètre retenu et les exclusions soient justifiés dans le rapport de
vérification.

Le niveau visé par ce projet et la justification des exclusions sont présentés
au chapitre~\ref{chap:modelisation}.

\subsection{Catalogues de risques}

Les catalogues de risques suivants sont pris en compte, conformément au sujet :

\begin{itemize}
  \item \textbf{OWASP Top 10:2025}~\cite{owasptop10} pour les risques
  applicatifs web ;
  \item \textbf{OWASP API Security Top 10}~\cite{owaspapi} pour les risques
  propres aux API REST, avec une attention particulière à l'authentification, à
  l'autorisation, au contrôle d'accès, à la validation des entrées, à
  l'exposition des données et à la protection contre les abus ;
  \item \textbf{OWASP Top 10 for LLM Applications}~\cite{owaspllm} pour les
  risques liés aux fonctionnalités fondées sur un modèle de langage ;
  \item \textbf{OWASP Mobile Top 10}~\cite{owaspmobile} pour l'application
  mobile, dont le traitement est reporté à la seconde phase du projet.
\end{itemize}

% ===========================================================================
\chapter{Architecture de l'application}
\label{chap:archi-app}
% ===========================================================================

\begin{constat}{état de ce chapitre}
Ce chapitre documente les choix de conception arrêtés. L'implémentation
correspondante relève de la phase~2 et sera documentée au
chapitre~\ref{chap:dev-securise}. Les diagrammes définitifs seront produits en
phase~1 avec le modèle de menaces.
\end{constat}

\section{Choix du domaine fonctionnel}

Le cas d'étude retenu est une \textbf{place de marché d'articles vintage} :
des vendeurs publient des objets d'occasion uniques, des acheteurs les
acquièrent.

Ce domaine n'a pas été choisi pour son intérêt fonctionnel mais pour la
richesse de sa surface de sécurité. Le tableau~\ref{tab:surface} justifie
chaque fonctionnalité retenue par la catégorie de risque qu'elle permet de
démontrer.

\begin{table}[H]
\centering
\caption{Justification des fonctionnalités par la surface de sécurité créée}
\label{tab:surface}
\small
\begin{tabular}{@{}p{3.9cm}p{4.0cm}p{6.0cm}@{}}
\toprule
\textbf{Fonctionnalité} & \textbf{Surface} & \textbf{Risque démontré} \\
\midrule
Inscription, connexion & Authentification
& Authentification défaillante, bourrage d'identifiants, stockage faible des
mots de passe \\
\addlinespace
Rôles acheteur, vendeur, administrateur & Autorisation
& Élévation de privilèges, contrôle d'accès au niveau des fonctions \\
\addlinespace
Gestion par le vendeur de ses seules annonces & Propriété des ressources
& \textbf{IDOR / BOLA} --- la démonstration centrale du projet \\
\addlinespace
Téléversement d'images & Fichiers non fiables
& Téléversement malveillant, traversée de chemin, exécution côté serveur \\
\addlinespace
Recherche d'annonces & Entrée utilisateur en requête
& Injection SQL, énumération, moissonnage \\
\addlinespace
Commande, paiement simulé & Logique métier à valeur
& Abus de logique métier, \textbf{double vente d'un objet unique} \\
\addlinespace
Avis sur commande & Contenu stocké
& XSS stocké \\
\addlinespace
Modération & Interface privilégiée
& Contrôle d'accès sur fonction d'administration \\
\addlinespace
Génération de description par IA & Intégration d'un LLM
& Injection d'invite, traitement non sûr de la sortie, abus de consommation \\
\bottomrule
\end{tabular}
\end{table}

\begin{decision}{périmètre fonctionnel gelé}
Sont explicitement hors périmètre : le traitement réel des paiements, l'envoi de
courriels, l'authentification fédérée et la messagerie temps réel.

Aucun de ces éléments n'apporte une catégorie de risque que le périmètre retenu
ne couvre déjà, et chacun consommerait du temps appartenant à la pipeline.
Toute fonctionnalité ajoutée devra en remplacer une autre.
\end{decision}

\section{Composants et choix techniques}

\begin{table}[H]
\centering
\caption{Composants de l'application}
\label{tab:composants}
\small
\begin{tabular}{@{}llp{6.4cm}@{}}
\toprule
\textbf{Composant} & \textbf{Technologie} & \textbf{Justification} \\
\midrule
API métier & Express.js, TypeScript
& Écosystème imposé par le sujet (Projet 3) ; typage statique \\
\addlinespace
Base de données & PostgreSQL
& Transactions et verrouillage de ligne, nécessaires au contrôle de double
vente \\
\addlinespace
Accès aux données & Prisma
& Requêtes paramétrées par défaut, \asvsc{1.2.4} \\
\addlinespace
Validation & Zod
& Validation par schéma en couche de service de confiance, \asvsc{2.2.2} \\
\addlinespace
Mots de passe & Argon2id
& Fonction de dérivation approuvée, \asvsc{11.4.2} \\
\addlinespace
En-têtes HTTP & Helmet
& En-têtes de sécurité du navigateur, \asvsc{3.4.x} \\
\addlinespace
Front-end web & React, TypeScript, Vite
& Partage du contrat d'API et des types avec le back-end \\
\addlinespace
Client mobile & React Native (Expo)
& \emph{Reporté en phase 2 du projet} ; réutilise la même API \\
\addlinespace
Fournisseur LLM & Google Gemini (offre gratuite)
& Aucun coût ; quota contraint, ce qui rend le risque \llm{10} concret \\
\bottomrule
\end{tabular}
\end{table}

\captureaplacer{figures/architecture-composants.png}{Diagramme de composants de
l'application : front-end web, client mobile (phase 2), API REST, base de
données, stockage d'objets, et fournisseur LLM identifié comme frontière de
confiance externe. À produire en phase 1.}

\section{Interface de programmation REST}

Le tableau~\ref{tab:api} présente la surface d'API prévue. Deux points
d'entrée sont désignés dès la conception comme cibles des démonstrations de
vulnérabilité du chapitre~\ref{chap:vulnerabilites}.

\begin{table}[H]
\centering
\caption{Surface d'API REST}
\label{tab:api}
\small
\begin{tabular}{@{}llp{5.9cm}@{}}
\toprule
\textbf{Point d'entrée} & \textbf{Accès} & \textbf{Contrôle principal} \\
\midrule
\texttt{POST /api/auth/register}   & public & Politique de mot de passe, limitation de débit \\
\texttt{POST /api/auth/login}      & public & Message d'erreur générique, anti-automatisation \\
\texttt{POST /api/auth/refresh}    & cookie & Rotation et détection de réutilisation \\
\texttt{POST /api/auth/logout}     & authentifié & Invalidation côté serveur \\
\texttt{GET /api/listings}         & public & Requête paramétrée, pagination \\
\texttt{GET /api/listings/:id}     & public & Champs publics uniquement \\
\texttt{POST /api/listings}        & vendeur & Validation, garde d'affectation de masse \\
\texttt{PUT /api/listings/:id}     & \textbf{propriétaire} & \textbf{Cible de la démonstration V5} \\
\texttt{DELETE /api/listings/:id}  & propriétaire & Propriété revérifiée côté serveur \\
\texttt{POST /api/listings/:id/images} & propriétaire & Octets magiques, ré-encodage \\
\texttt{POST /api/orders}          & acheteur & Verrou transactionnel \\
\texttt{GET /api/orders}           & \textbf{propriétaire} & \textbf{Cible de démonstration BOLA} \\
\texttt{POST /api/reviews}         & acheteur & Assainissement en sortie \\
\texttt{GET /api/admin/listings}   & administrateur & Contrôle d'accès par fonction \\
\texttt{POST /api/ai/describe}     & vendeur & Gardes d'injection, plafond de jetons \\
\bottomrule
\end{tabular}
\end{table}

\section{Modèle de données}

Les entités prévues sont : \texttt{User}, \texttt{Listing},
\texttt{ListingImage}, \texttt{Order}, \texttt{Review} et
\texttt{RefreshToken}.

Une contrainte de conception mérite d'être soulignée : l'entité
\texttt{Listing} doit porter un état permettant de distinguer un article
disponible d'un article réservé ou vendu. Cet état est le support du contrôle de
non-duplication de vente (\asvsc{2.3.4}), lequel constitue la particularité
d'une place de marché d'objets uniques par rapport à un commerce de biens en
stock.

\section{Classification des données}

\begin{constat}{à compléter en phase 1}
La classification des données et les exigences de protection associées à chaque
niveau constituent une exigence documentaire de l'ASVS (\asvsc{14.1.1} et
\asvsc{14.1.2}). Elles seront produites en phase~1, en même temps que le modèle
de menaces, et reportées ici.
\end{constat}

\section{Architecture indifférente au client}

L'API est conçue dès l'origine pour servir indifféremment le front-end web et
le futur client mobile. Ce choix a une conséquence pratique sur le report de
l'application mobile exposé au chapitre~\ref{chap:introduction} : l'ensemble des
contrôles côté serveur, le modèle de menaces, la pipeline, la suite de tests
d'abus et les exigences ASVS côté serveur sont réutilisés sans modification.

Les travaux propres au mobile se limitent au client lui-même et aux risques qui
lui sont spécifiques --- stockage local, épinglage de certificat, intégrité de
l'artefact --- recensés au chapitre~\ref{chap:limites}.

% ===========================================================================
\chapter{Architecture de la pipeline DevSecOps}
\label{chap:archi-pipeline}
% ===========================================================================

\begin{constat}{état de ce chapitre}
Ce chapitre présente la conception d'ensemble de la pipeline et la
correspondance entre pratiques, outils et étages. Les étages de la phase~0 sont
en service ; les étages d'analyse de sécurité sont mis en œuvre en phase~5 et
documentés au chapitre~\ref{chap:stages}.
\end{constat}

\section{Principe de conception}

La pipeline est structurée par la cartographie des pratiques sur les phases du
cycle de vie DevSecOps établie par Prates et Pereira~\cite{prates2025}. Cette
cartographie répond à la question de l'emplacement : elle indique à quel moment
de la chaîne chaque pratique doit intervenir.

\begin{table}[H]
\centering
\caption{Correspondance entre phases du cycle DevSecOps, pratiques et étages de
notre pipeline}
\label{tab:cycle-pipeline}
\small
\begin{tabular}{@{}llp{7.0cm}@{}}
\toprule
\textbf{Phase du cycle} & \textbf{Pratiques} & \textbf{Mise en œuvre} \\
\midrule
Create
& \pid{P8} \pid{P9} \pid{P20}
& Crochets \emph{pre-commit}, \texttt{CODEOWNERS}, guide de codage sécurisé \\
\addlinespace
Plan
& \pid{P7} \pid{P23} \pid{P24}
& Modélisation des menaces, analyse de risques, périmètre ASVS \\
\addlinespace
Adapt
& \pid{P6} \pid{P12}
& Intégration continue, infrastructure en tant que code \\
\addlinespace
Verify
& \pid{P1} \pid{P2} \pid{P3} \pid{P4} \pid{P11}
& SAST, DAST, suite de tests d'abus, analyse de vulnérabilités, revue \\
\addlinespace
Preproduction
& \pid{P4} \pid{P16} \pid{P22}
& Analyse d'image, test d'intrusion, durcissement \\
\addlinespace
Release
& \pid{P10}
& Déploiement continu conditionné à la réussite des contrôles \\
\addlinespace
Predict / Detect
& \pid{P5} \pid{P18}
& Supervision continue, journalisation, alertes \\
\addlinespace
Respond
& \pid{P28}
& Procédure de gestion des incidents \\
\addlinespace
Configure
& \pid{P30} \pid{P21}
& Gestion des secrets, gestion de configuration \\
\bottomrule
\end{tabular}
\end{table}

\captureaplacer{figures/architecture-pipeline.png}{Schéma de la pipeline
DevSecOps : étages, déclencheurs et points de blocage, superposés aux phases du
cycle de vie. À produire en phase 5, lorsque tous les étages seront en service.}

\section{Étages prévus et outils retenus}

\subsection{Règle de traçabilité des outils}

Un principe a été appliqué à l'ensemble des choix d'outillage.

\begin{decision}{outils issus de la source scientifique}
À l'exception de deux outils explicitement signalés, tous les outils retenus
figurent dans le tableau~8 de Prates et Pereira~\cite{prates2025}. Chaque choix
est donc rattachable à la source de référence du module.

Les écarts sont signalés comme tels et justifiés en propre. Cette distinction
est maintenue dans tout le rapport : elle évite de présenter une préférence
personnelle comme une recommandation de la littérature.
\end{decision}

\subsection{Tableau des étages}

\begin{table}[H]
\centering
\caption{Étages de sécurité de la pipeline}
\label{tab:etages}
\small
\begin{tabular}{@{}llp{3.5cm}p{4.3cm}@{}}
\toprule
\textbf{Étage} & \textbf{Prat.} & \textbf{Outil} & \textbf{Déclencheur} \\
\midrule
Détection de secrets & \pid{P30} & Gitleaks & chaque envoi, historique complet \\
Hygiène du dépôt     & \pid{P8}  & pre-commit & chaque envoi \\
SAST                 & \pid{P1}  & SonarQube, ESLint-security & envoi et \emph{pull request} \\
Composition logicielle & \pid{P13} & OWASP Dependency-Check & \emph{pull request} et nocturne \\
Dépendances JS       & \pid{P4}  & Retire.js & \emph{pull request} \\
Image de conteneur   & \pid{P4}  & Trivy & après construction d'image \\
Infrastructure       & \pid{P4}  & Checkov & modification de manifeste \\
Tests d'abus         & \pid{P3}  & Jest, Supertest & \emph{pull request} \\
DAST                 & \pid{P2}  & OWASP ZAP & référence en PR, complet en nocturne \\
Conformité continue  & \pid{P17} & Checkov, politique & \emph{pull request} \\
\bottomrule
\end{tabular}
\end{table}

Retire.js mérite un commentaire particulier : cet outil est spécifique à
l'écosystème JavaScript, ce qui constitue une adéquation directe au périmètre du
Projet~3.

\section{Politique de blocage}

\begin{constat}{la politique de blocage est elle-même une décision de sécurité}
Déterminer quels résultats font échouer la construction et lesquels ne
produisent qu'un avertissement n'est pas un réglage technique mineur.

Faire échouer la construction sur chaque résultat de faible gravité paralyse le
travail et enseigne à l'équipe le contournement de la pipeline. Ne rien faire
échouer rend la pipeline décorative. La position intermédiaire doit être
justifiée explicitement, et non supposée.
\end{constat}

La politique retenue sera consignée dans le registre ADR~0008 en phase~5, et
reproduite au chapitre~\ref{chap:stages}. Elle précisera pour chaque outil le
seuil de gravité bloquant, ainsi que la procédure d'acceptation d'un résultat
--- par qui, sur quel motif, et avec quelle échéance de réexamen.

\section{Séparation entre contrôles rapides et analyses profondes}

Deux contraintes s'opposent. Les contrôles doivent être suffisamment rapides
pour s'exécuter à chaque \emph{pull request} sans entraver le travail ; les
analyses dynamiques complètes ne peuvent pas l'être.

La conception retenue sépare donc les déclencheurs : analyse ZAP de référence et
analyse de composition rapide sur chaque \emph{pull request}, analyse ZAP
complète et analyse de composition approfondie sur déclenchement nocturne
planifié.

Ce choix est à la fois une décision de sécurité et une décision de performance.
Son effet mesuré est présenté au chapitre~\ref{chap:performance}.

% ===========================================================================
\chapter{Phase 0 --- Fondations et gouvernance de sécurité}
\label{chap:phase0}
% ===========================================================================

\section{Stratégie : la sécurité avant le code}

\subsection{Justification de l'ordonnancement}

La première phase de réalisation ne produit aucune ligne de code applicatif.
Elle met en place la gouvernance de sécurité du dépôt et la pipeline
d'intégration continue, de sorte que \emph{toute} contribution ultérieure soit
contrôlée dès son premier envoi.

Ce choix découle directement de deux principes énoncés au
chapitre~\ref{chap:fondements}.

Le premier est la pratique \pid{P8} \emph{Shift left security}, définie comme
l'intégration des considérations de sécurité le plus tôt possible dans le cycle
de développement. L'étape la plus précoce techniquement atteignable n'est pas
l'intégration continue : c'est le poste du développeur, avant même la création
du commit.

Le second est l'avertissement formulé par l'ASVS~5.0, selon lequel concevoir la
cryptographie au fur et à mesure ou la moderniser après coup coûtera
inévitablement beaucoup plus cher à mettre en œuvre de manière sécurisée que de
l'intégrer dès le départ~\cite{owaspasvs5}. Cette observation se généralise à
l'ensemble des contrôles.

\begin{decision}{pipeline avant code applicatif}
La pipeline d'intégration continue est créée et rendue fonctionnelle
\emph{avant} l'existence du code applicatif.

La conséquence technique est que la pipeline doit réussir sur un projet vide.
Le travail consistant à satisfaire cette contrainte est décrit à la
section~\ref{sec:ci-projet-vide}. Le bénéfice est qu'il devient impossible, par
construction, d'ajouter du code non contrôlé : la règle de protection de branche
dispose d'un nom de contrôle stable dès le premier jour.
\end{decision}

\subsection{Pratiques couvertes par cette phase}

Six pratiques sont mises en œuvre dans cette phase, dont quatre de niveau
Tier~1, sans qu'aucune logique métier n'ait été écrite.

\begin{table}[H]
\centering
\caption{Pratiques couvertes par la phase 0}
\label{tab:phase0-pratiques}
\small
\begin{tabular}{@{}llp{7.4cm}@{}}
\toprule
\textbf{Id.} & \textbf{Niveau} & \textbf{Mise en œuvre} \\
\midrule
\pid{P6}  & Tier 1 & Pipeline GitHub Actions à trois tâches \\
\pid{P8}  & Tier 1 & Crochets \emph{pre-commit} avec Gitleaks \\
\pid{P9}  & Tier 1 & \texttt{CODEOWNERS}, modèle de \emph{pull request} \\
\pid{P11} & Tier 1 & Protection de branche, revue obligatoire \\
\pid{P15} & Tier 2 & Dépôt Git, historique tracé \\
\pid{P30} & Tier 2 & Configuration Gitleaks, règles spécifiques au projet \\
\bottomrule
\end{tabular}
\end{table}

\section{Gouvernance du dépôt}

\subsection{Organisation de l'équipe}

Le dépôt est hébergé sur GitHub à l'adresse
\url{https://github.com/edothecreator/vintage-marketplace-devsecops}.

Les quatre membres du projet y disposent d'un accès en écriture. La répartition
des domaines de responsabilité, formalisée dans le fichier \texttt{CODEOWNERS}
(section~\ref{sec:codeowners}), est la suivante : pipeline et sécurité ;
back-end et API ; front-end et suite de tests d'abus ; infrastructure,
supervision et documentation.

\subsection{Arborescence}

L'arborescence a été définie dès cette phase afin que les artefacts de preuve
disposent d'un emplacement déterminé avant d'être produits.

\begin{lstlisting}[language=bash,caption={Arborescence du dépôt}]
backend/            API REST Express.js + TypeScript
frontend/           front-end web
docs/
  adr/              registres de decisions d'architecture
  threat-model/     modele de menaces et diagrammes
  phases/           suivi des taches par phase
evidence/           preuves par pratique (P1..P11, tier2)
tasks/              fiches de tache pour les membres de l'equipe
.github/workflows/  definitions de la pipeline
\end{lstlisting}

Le répertoire \texttt{evidence/} comporte un sous-répertoire par pratique de
niveau Tier~1. Cette organisation répond à une difficulté pratique identifiée
dès le départ : les journaux d'exécution de pipeline expirent, les branches sont
supprimées, les tableaux de bord sont réinitialisés. Une preuve non capturée au
moment où elle est produite est généralement irrécupérable.

\begin{decision}{capture des preuves intégrée à la définition de « terminé »}
La capture des preuves fait partie de la définition de l'achèvement de chaque
tâche, et non d'une tâche distincte planifiée en fin de projet.
\end{decision}

\section{P6 --- Intégration continue}
\label{sec:ci}

\subsection{Choix de la plate-forme}

L'étude de Prates et Pereira recense cinq outils d'intégration continue :
Jenkins (15 mentions), GitLab CI (8), GitHub Actions (6), Travis CI (3) et
Circle CI (2)~\cite{prates2025}. Plusieurs options étaient donc justifiables au
regard de la source.

\begin{decision}{GitHub Actions (registre ADR 0002)}
GitHub Actions a été retenu pour les motifs suivants :

\begin{itemize}
  \item la plate-forme est native à l'hébergement du dépôt et des
  collaborateurs, ce qui permet l'intégration directe des contrôles obligatoires
  aux \emph{pull requests} --- condition nécessaire pour que \pid{P11} soit
  contraignant et non indicatif ;
  \item aucune infrastructure à maintenir, à l'inverse de Jenkins dont
  l'hébergement, la maintenance et la sécurisation auraient constitué un projet
  en soi, au détriment du temps consacré à l'ingénierie de sécurité ;
  \item l'outil Dependabot est disponible gratuitement, ce qui apporte la
  pratique Tier~2 \pid{P27} \emph{Patch management} à coût quasi nul.
\end{itemize}

GitLab CI, deuxième outil le plus cité, a été écarté car le dépôt et les
collaborateurs sont sur GitHub, et parce que ses fonctionnalités de sécurité les
plus attrayantes relèvent d'offres payantes : nous aurions reconstitué les mêmes
analyses avec les mêmes outils libres, après avoir changé de plate-forme sans
bénéfice.
\end{decision}

\subsection{Structure de la pipeline}

La pipeline de la phase 0 comporte trois tâches.

\begin{table}[H]
\centering
\caption{Tâches de la pipeline en phase 0}
\label{tab:ci-jobs}
\small
\begin{tabular}{@{}llp{6.8cm}@{}}
\toprule
\textbf{Tâche} & \textbf{Pratique} & \textbf{Rôle} \\
\midrule
\texttt{Secret scan (gitleaks)}
& \pid{P30}
& Analyse de l'\emph{historique complet} du dépôt, et non du seul diff \\
\addlinespace
\texttt{Repository hygiene (pre-commit)}
& \pid{P8}
& Exécute en CI les mêmes crochets que localement \\
\addlinespace
\texttt{Build and test}
& \pid{P6}
& Construction et tests ; conditionnelle à l'existence du code \\
\bottomrule
\end{tabular}
\end{table}

Deux choix de configuration méritent d'être explicités.

Le premier concerne les \textbf{permissions du jeton d'exécution}, réduites à la
seule lecture du contenu par défaut. Toute tâche nécessitant davantage doit le
demander explicitement. Il s'agit de l'application du principe de moindre
privilège à la chaîne de production elle-même.

Le second concerne la \textbf{tâche d'hygiène}, qui exécute en intégration
continue exactement les mêmes crochets que ceux installés sur le poste du
développeur. Cette redondance est volontaire : elle garantit que l'option
\texttt{-{}-no-verify}, qui permet localement de contourner un crochet, ne
permet pas de contourner durablement le contrôle, car l'envoi échoue malgré
tout.

\begin{lstlisting}[language=yamlperso,caption={Extrait de
\texttt{.github/workflows/ci.yml} --- principe de moindre privilège et
contrôle de concurrence}]
name: CI

on:
  push:
    branches: ["main", "develop", "feat/**", "fix/**", "demo/**"]
  pull_request:
    branches: ["main", "develop"]

# Moindre privilege par defaut. Toute tache necessitant davantage
# doit le demander explicitement.
permissions:
  contents: read

# Un nouvel envoi sur la meme reference annule l'execution precedente,
# sauf sur main.
concurrency:
  group: ${{ github.workflow }}-${{ github.ref }}
  cancel-in-progress: ${{ github.ref != 'refs/heads/main' }}
\end{lstlisting}

\subsection{Réussir sur un projet vide}
\label{sec:ci-projet-vide}

La contrainte énoncée en début de chapitre --- rendre la pipeline fonctionnelle
avant l'existence du code --- impose que la tâche de construction réussisse sur
un dépôt sans application.

La solution retenue consiste à détecter la présence du code applicatif et à
n'exécuter les étapes Node.js que conditionnellement.

\begin{lstlisting}[language=yamlperso,caption={Détection conditionnelle du code
applicatif}]
- name: Detect application scaffold
  id: detect
  run: |
    if [ -f backend/package.json ]; then
      echo "scaffold=true" >> "$GITHUB_OUTPUT"
    else
      echo "scaffold=false" >> "$GITHUB_OUTPUT"
    fi

# ... etapes Node.js conditionnees par steps.detect.outputs.scaffold

- name: No scaffold yet
  if: steps.detect.outputs.scaffold == 'false'
  run: |
    echo "backend/package.json non present - la phase 2 n'a pas commence."
    echo "Cette tache reussit volontairement afin que la protection de"
    echo "branche dispose d'un controle requis stable des la phase 0."
\end{lstlisting}

\subsection{Résultat}

\begin{verification}{pipeline fonctionnelle}
La pipeline a été exécutée avec succès. Les trois tâches aboutissent, et les
étapes Node.js sont bien ignorées en l'absence de code applicatif, ce qui valide
la logique conditionnelle.
\end{verification}

\captureprete{figures/p06-pipeline-run.png}{Première exécution de la pipeline :
les trois tâches réussissent. Les étapes Node.js sont ignorées, la tâche
« No scaffold yet » confirmant la détection correcte de l'absence de code
applicatif.}

\section{P8 --- Shift left security}
\label{sec:shift-left}

\subsection{Choix du mécanisme de crochets}

\begin{decision}{cadre \emph{pre-commit} (registre ADR 0003)}
Le cadre \texttt{pre-commit}~\cite{precommit} a été retenu plutôt que Husky,
pourtant plus idiomatique dans un projet Node.

Le motif est opérationnel : Gitleaks n'est pas un paquet npm. Avec Husky, il
aurait fallu soit demander à chaque membre de l'équipe d'installer manuellement
le binaire Gitleaks, soit invoquer Docker à chaque commit --- lent, et
inopérant pour quiconque n'a pas Docker en fonctionnement. Le cadre
\texttt{pre-commit} télécharge et épingle lui-même la version du binaire : une
seule commande suffit pour toute l'équipe.
\end{decision}

\subsection{Composition des crochets}

Les crochets installés sont délibérément rapides.

\begin{decision}{rapidité des crochets locaux}
Les analyses lentes --- SAST, SCA, DAST --- restent en intégration continue et
ne sont pas exécutées à chaque commit.

Un crochet lent enseigne aux développeurs l'usage de \texttt{-{}-no-verify}, ce
qui annule le contrôle. La rapidité n'est donc pas ici un critère de confort
mais une condition d'efficacité du contrôle.
\end{decision}

Les crochets retenus sont : Gitleaks (détection de secrets), refus des clés
privées, refus des fichiers volumineux, validation syntaxique YAML et JSON,
normalisation des fins de ligne, et un crochet local refusant explicitement la
mise en index d'un fichier \texttt{.env} réel.

\subsection{Omission volontaire}

Les crochets ESLint, Prettier et de vérification de types n'ont
\emph{pas} été configurés à cette phase.

\begin{decision}{report des crochets de qualité à la phase 2}
Sur un projet vide, ces crochets échoueraient immédiatement. La première
expérience de l'équipe avec l'outillage de sécurité serait alors un chemin de
commit cassé --- le moyen le plus rapide d'enseigner le contournement.

Ils seront ajoutés en phase 2, dès l'existence de \texttt{backend/package.json}.
\end{decision}

\section{P30 --- Gestion des secrets}
\label{sec:secrets}

\subsection{Règles spécifiques au projet}

La configuration Gitleaks part du jeu de règles par défaut maintenu en amont,
puis l'étend de quatre règles propres au projet.

\begin{lstlisting}[caption={Extrait de \texttt{.gitleaks.toml} --- règles
spécifiques au projet}]
[extend]
useDefault = true

[[rules]]
id = "gemini-api-key"
description = "Cle API Google Gemini / AI Studio"
regex = '''AIza[0-9A-Za-z\-_]{35}'''
tags = ["key", "google", "gemini"]

[[rules]]
id = "generic-jwt-secret-assignment"
description = "Secret de signature JWT code en dur"
regex = '''(?i)(jwt[_-]?secret|token[_-]?secret|signing[_-]?key)\s*[:=]\s*['"][^'"\s]{8,}['"]'''
tags = ["secret", "jwt"]

[[rules]]
id = "hardcoded-db-url"
description = "Chaine de connexion avec identifiants en ligne"
regex = '''(?i)(postgres|postgresql|mysql|mongodb(\+srv)?)://[^:\s'"]+:[^@\s'"]+@'''
tags = ["credential", "database"]
\end{lstlisting}

La liste d'autorisation est volontairement restreinte aux valeurs de
substitution, aux fixtures de test et à la documentation. Une liste
d'autorisation trop large produirait une configuration qui semble fonctionner
tout en ne détectant rien.

\subsection{Vérification de l'efficacité des règles}

Une configuration de détection non testée n'est pas un contrôle. Les règles ont
donc été vérifiées par introduction volontaire d'un fichier de test contenant
une fausse clé Gemini et un secret de signature JWT codé en dur.

\begin{verification}{les règles de détection fonctionnent}
Les deux règles ont déclenché. La sortie de l'outil est reproduite ci-dessous,
les valeurs sensibles étant automatiquement masquées par l'option
\texttt{-{}-redact} :

\begin{lstlisting}[basicstyle=\ttfamily\tiny,numbers=none]
gitleaks (block secrets before commit)...................................Failed

Finding:     geminiApiKey: "REDACTED0a",
Secret:      REDACTED
RuleID:      gemini-api-key
Entropy:     5.285402
Tags:        [key google gemini]
File:        backend/src/config/TEMP-leak-test.ts
Line:        3

Finding:     REDACTED: "sup3rS3cr3tSigningKeyForTokens"
Secret:      REDACTED
RuleID:      generic-jwt-secret-assignment
Entropy:     2.725481
Tags:        [secret jwt]
File:        backend/src/config/TEMP-leak-test.ts
\end{lstlisting}

Le fichier de test a été supprimé immédiatement après vérification. Cette
vérification présente un second intérêt : elle valide par avance la
démonstration V2 du chapitre~\ref{chap:vulnerabilites}.
\end{verification}

\section{P9 --- Collaboration}
\label{sec:codeowners}

\subsection{Attribution de la propriété du code}

Le fichier \texttt{CODEOWNERS} associe un ou plusieurs propriétaires à chaque
chemin du dépôt. Combiné à l'option de protection de branche exigeant la revue
des propriétaires, il transforme la revue en exigence structurelle.

Les chemins sensibles à la sécurité --- définitions de pipeline, configuration
Gitleaks, crochets, politique de sécurité --- sont attribués au responsable de
la pipeline, car une modification de ces fichiers peut désactiver silencieusement
un contrôle. Les répertoires d'authentification et d'intergiciels exigent deux
relecteurs.

\begin{constat}{\texttt{CODEOWNERS} échoue silencieusement}
GitHub ignore sans avertissement les noms d'utilisateur inexistants ou ne
disposant pas d'un accès en écriture. Le fichier peut donc avoir toutes les
apparences d'un contrôle fonctionnel sans rien imposer.

Les quatre identifiants ont par conséquent été validés par appel à l'API avant
usage. La validité du fichier est contrôlable par le point d'accès dédié :

\begin{lstlisting}[language=bash,numbers=none]
gh api repos/<owner>/<repo>/codeowners/errors
\end{lstlisting}
\end{constat}

\subsection{Modèle de \emph{pull request}}

Un modèle de \emph{pull request} impose une liste de contrôle de sécurité à
chaque contribution, couvrant la validation des entrées, le contrôle d'accès,
les secrets, les dépendances, la journalisation, et une section spécifique au
code produit avec l'assistance d'une IA générative, conformément à la
section~7 du sujet.

La consigne retenue est qu'une case non cochée doit être justifiée par écrit :
une case laissée vide sans explication est traitée comme non revue.

\captureaplacer{figures/p09-collaborators.png}{Page des collaborateurs du dépôt
(Settings $\rightarrow$ Collaborators) montrant les quatre membres de l'équipe
avec accès en écriture. À capturer après acceptation des invitations.}

\section{P11 --- Revue de code}
\label{sec:revue-code}

\subsection{Protection de la branche principale}

La revue de code n'est une pratique mise en œuvre que si elle est contraignante.
La configuration appliquée à la branche \texttt{main} est la suivante.

\begin{table}[H]
\centering
\caption{Règles de protection appliquées à la branche \texttt{main}}
\label{tab:branch-protection}
\small
\begin{tabular}{@{}p{8.2cm}p{5.4cm}@{}}
\toprule
\textbf{Règle} & \textbf{Valeur} \\
\midrule
Revues approbatrices requises & 1 \\
Revue des propriétaires du code requise & oui \\
Invalidation des revues obsolètes après envoi & oui \\
Approbation du dernier envoi requise & oui \\
Contrôles d'état requis & 3, en mode strict \\
Résolution des discussions requise & oui \\
Envois forcés & interdits \\
Suppression de branche & interdite \\
\bottomrule
\end{tabular}
\end{table}

\begin{constat}{les noms des contrôles requis doivent correspondre exactement}
Les trois contrôles d'état requis sont désignés par le \emph{nom} des tâches de
la pipeline. Un renommage dans \texttt{ci.yml} détache silencieusement la règle
de protection, qui continue d'exister tout en n'exigeant plus rien.

La correspondance exacte entre les noms de tâches effectivement exécutés et les
noms exigés par la protection a donc été vérifiée explicitement.
\end{constat}

\captureaplacer{figures/p11-branch-protection-reviews.png}{Règles de protection
de la branche \texttt{main} (Settings $\rightarrow$ Branches $\rightarrow$ Edit) :
exigence de \emph{pull request}, une revue approbatrice, invalidation des revues
obsolètes, revue des propriétaires du code, approbation du dernier envoi.}

\captureaplacer{figures/p11-branch-protection-status-checks.png}{Suite des règles
de protection : les trois contrôles d'état requis, l'exigence de branche à jour,
et l'interdiction des envois forcés et des suppressions.}

\section{Difficultés rencontrées et décisions induites}

Trois obstacles techniques ont été rencontrés durant cette phase. Ils sont
rapportés ici parce qu'ils ont chacun entraîné une décision documentée, et parce
que leur résolution fait partie du travail d'ingénierie accompli.

\subsection{Protection de branche indisponible sur dépôt privé}

Le dépôt a initialement été créé en mode privé, choix instinctivement plus
prudent. Les tentatives d'activation de la protection de branche, puis des
\emph{rulesets}, ont toutes deux échoué :

\begin{lstlisting}[numbers=none,basicstyle=\ttfamily\scriptsize]
HTTP 403 - Upgrade to GitHub Pro or make this repository public
            to enable this feature.
\end{lstlisting}

L'offre GitHub Free n'autorise ni les branches protégées ni les \emph{rulesets}
sur un dépôt privé. Trois options s'ouvraient : rendre le dépôt public, obtenir
GitHub Pro gratuitement via le \emph{Student Developer Pack}, ou renoncer à la
protection de branche.

La troisième option a été écartée d'emblée : elle aurait réduit la preuve de
deux pratiques de niveau Tier~1 de « la fusion est techniquement bloquée » à
« nous avons rédigé un modèle », ce qui ne démontre pas un contrôle mis en
œuvre.

\begin{decision}{dépôt public (registre ADR 0004)}
Le dépôt a été rendu public. Outre le déblocage de la protection de branche,
ce choix apporte un bénéfice indépendant et matériel : \textbf{les minutes
d'exécution GitHub Actions ne sont pas comptabilisées sur un dépôt public}.

Ce point n'est pas accessoire. Les analyses ZAP complètes et SonarQube sont les
étages les plus coûteux en temps, et le chapitre~\ref{chap:vulnerabilites}
déclenche à lui seul sept exécutions complètes de la pipeline sur sept branches
de démonstration. Sur un dépôt privé, l'allocation gratuite de 2000 minutes
mensuelles nous aurait contraints à rationner les analyses pour des motifs
étrangers à l'ingénierie de sécurité.

La contrepartie est que le code et la configuration sont lisibles publiquement,
ce qui est acceptable pour un travail académique sans contenu propriétaire, et
qu'\emph{aucun} secret réel ne doit jamais entrer dans le dépôt --- ce qui
était déjà la politique, outillée par les crochets locaux et l'analyse
d'historique complet en CI.
\end{decision}

\subsection{Interférence avec la démonstration V2}

Le passage en dépôt public active la protection contre l'envoi de secrets de
GitHub, qui analyse les motifs propres à des fournisseurs identifiés et peut
rejeter l'envoi.

\begin{constat}{un contrôle externe peut masquer le nôtre}
La démonstration V2 prévue au chapitre~\ref{chap:vulnerabilites} introduit
volontairement un secret codé en dur afin que Gitleaks le détecte et fasse
échouer la construction.

Si ce secret est une clé API Google de la forme \texttt{AIza...}, GitHub rejette
l'envoi \emph{avant} que notre pipeline ne s'exécute. Notre contrôle paraîtrait
alors avoir échoué, alors qu'il n'aurait simplement jamais été sollicité.

La démonstration V2 utilise donc le motif de secret de signature JWT, détecté
par notre règle \texttt{generic-jwt-secret-assignment} et non reconnu par
l'analyseur de GitHub. Les deux règles ont été vérifiées
(section~\ref{sec:secrets}).

Cette substitution améliore d'ailleurs la démonstration : un secret de signature
codé en dur est une erreur de développement réaliste, et sa détection repose sur
\emph{notre} configuration plutôt que sur la liste de motifs d'un tiers.
\end{constat}

\subsection{Blocage du gestionnaire d'identifiants Git}

Les opérations Git distantes se sont initialement interrompues sans message.
Le diagnostic a établi que le gestionnaire d'identifiants configuré
(\emph{Git Credential Manager}) tentait d'afficher une boîte de dialogue
d'authentification impossible à présenter dans le contexte d'exécution, et
attendait indéfiniment une saisie.

La résolution a consisté à router les identifiants Git par le jeton de l'outil
en ligne de commande GitHub :

\begin{lstlisting}[language=bash,numbers=none]
gh auth setup-git
\end{lstlisting}

Ce point est consigné car les autres membres de l'équipe rencontreront le même
obstacle au premier clonage. Il figure à ce titre dans la documentation
d'installation.

\section{Registres de décisions d'architecture}

Quatre registres de décisions ont été produits durant cette phase.

\begin{table}[H]
\centering
\caption{Registres de décisions d'architecture de la phase 0}
\label{tab:adr-phase0}
\small
\begin{tabular}{@{}llp{8.2cm}@{}}
\toprule
\textbf{Réf.} & \textbf{Statut} & \textbf{Objet} \\
\midrule
ADR 0001 & Accepté & Tenue de registres de décisions d'architecture \\
ADR 0002 & Accepté & GitHub Actions comme plate-forme CI/CD \\
ADR 0003 & Accepté & Cadre \texttt{pre-commit} pour les crochets \emph{shift left} \\
ADR 0004 & Accepté & Visibilité publique du dépôt et atténuation pour V2 \\
\bottomrule
\end{tabular}
\end{table}

Le registre ADR 0001 justifie la pratique elle-même. Il s'appuie sur
l'observation que le raisonnement conduisant à une décision est ce qui s'efface
le plus vite, et sur l'exigence de l'ASVS~5.0 relative aux décisions de sécurité
documentées. Les registres sont immuables : une décision révisée donne lieu à un
nouveau registre qui remplace le précédent, l'original demeurant afin que
l'historique du raisonnement subsiste.

\section{Bilan de la phase}

\begin{table}[H]
\centering
\caption{État d'avancement à l'issue de la phase 0}
\label{tab:bilan-phase0}
\small
\begin{tabular}{@{}llp{6.6cm}@{}}
\toprule
\textbf{Élément} & \textbf{État} & \textbf{Observation} \\
\midrule
Dépôt, arborescence, configuration & \fait{} & 4 commits \\
Pipeline d'intégration continue    & \fait{} & 2 exécutions, toutes réussies \\
Crochets \emph{pre-commit}         & \fait{} & efficacité vérifiée \\
Règles de détection de secrets     & \fait{} & deux règles déclenchées en test \\
Protection de branche              & \fait{} & vérifiée par API \\
\texttt{CODEOWNERS}                & \encours{} & en attente d'acceptation des invitations \\
Registres de décisions             & \fait{} & ADR 0001 à 0004 \\
Captures de preuve                 & \encours{} & 1 fournie sur 3 \\
\bottomrule
\end{tabular}
\end{table}

L'ingénierie de la phase est achevée. Deux éléments de preuve restent à
collecter : les captures des réglages de plate-forme, qui ne peuvent être
produites que par un administrateur du dépôt et constituent la seule preuve
qu'un contrôle est effectivement activé, et l'acceptation des invitations par
les trois autres membres, préalable au fonctionnement effectif de l'exigence de
revue par les propriétaires du code.

Aucun de ces deux points ne bloque la phase suivante.


% --- Chapitres a venir, au fur et a mesure de l'avancement du projet -------
% ===========================================================================
\chapter{Modélisation des menaces et exigences de sécurité}
\label{chap:modelisation}
% ===========================================================================

\begin{constat}{chapitre à rédiger --- phase 1}
Ce chapitre documentera les livrables de la phase~1, qui ne produit
volontairement aucun code applicatif.

Contenu prévu :
\begin{itemize}
  \item \pid{P7} modélisation des menaces : analyse STRIDE par flux de données,
  réalisée avec OWASP Threat Dragon~\cite{threatdragon} ;
  \item table de traçabilité menace $\rightarrow$ mesure d'atténuation
  $\rightarrow$ exigence ASVS ;
  \item \pid{P23} analyse de risques : cotation vraisemblance $\times$ impact,
  priorisation des contrôles, risques explicitement acceptés ;
  \item \pid{P24} recueil des exigences de sécurité : niveau ASVS visé
  (niveau~1 complet et niveau~2 ciblé sur les chapitres V2, V4, V5, V6, V8, V13
  et V16), avec justification écrite des exclusions (V10 et V17) ;
  \item décisions de sécurité documentées exigées par l'ASVS~5.0 : règles de
  validation, règles d'autorisation, types de fichiers autorisés, durées de
  session, politique de mot de passe, niveaux de protection des données,
  inventaire de journalisation, calendrier de rotation des secrets ;
  \item classification des données et exigences de protection associées ;
  \item correspondance entre les risques du Top~10 OWASP pour les applications
  LLM et les contrôles prévus.
\end{itemize}

Registres de décisions attendus : ADR~0005 (niveau ASVS), ADR~0006 (décisions
de sécurité documentées), ADR~0007 (périmètre de l'augmentation par IA).
\end{constat}

% ===========================================================================
\chapter{Developpement securise du back-end}
\label{chap:dev-securise}
% ===========================================================================

\begin{constat}{chapitre a rediger --- phase 2}
P20 codage securise : controles implementes avec identifiant ASVS en commentaire ; stockage des mots de passe Argon2id ; validation par schema ; intergiciel central d'autorisation ; garde d'affectation de masse ; telechargement de fichiers avec verification des octets magiques ; verrouillage transactionnel contre la double vente ; limitation de debit ; en-tetes de securite ; gestion des erreurs ; journalisation de securite. Table de tracabilite exigence ASVS vers emplacement dans le code. Specification OpenAPI.
\end{constat}

% ===========================================================================
\chapter{Tests de securite automatises}
\label{chap:tests}
% ===========================================================================

\begin{constat}{chapitre a rediger --- phase 3}
P3 tests de securite automatises, troisieme pratique la plus citee (56 mentions). Suite de cas d'abus Jest et Supertest, chaque test nommant l'exigence ASVS qu'il verifie. Groupes : autorisation (IDOR, BOLA, elevation de privileges), logique metier (double vente concurrente), telechargement de fichiers, injection et traitement des sorties, authentification et jetons, affectation de masse, anti-automatisation. Table de correspondance test vers exigence ASVS.
\end{constat}

% ===========================================================================
\chapter{Etages de securite de la pipeline}
\label{chap:stages}
% ===========================================================================

\begin{constat}{chapitre a rediger --- phase 5}
P1 SAST avec SonarQube et ESLint-security ; P2 DAST avec OWASP ZAP ; P4 analyse de vulnerabilites avec Retire.js, Trivy et Checkov ; P13 analyse de composition avec OWASP Dependency-Check et generation de SBOM ; P17 conformite continue ; P27 gestion des correctifs avec Dependabot. Pour chaque etage : outil, declencheur, politique de blocage, verification que le blocage fonctionne effectivement. Registre ADR 0008 sur la politique de blocage.
\end{constat}

% ===========================================================================
\chapter{Conteneurisation et deploiement}
\label{chap:conteneurisation}
% ===========================================================================

\begin{constat}{chapitre a rediger --- phase 6}
P10 deploiement continu conditionne a la reussite de tous les controles ; P12 infrastructure en tant que code ; P22 durcissement de l'environnement. Dockerfile multi-etapes, execution sans privileges root, image de base minimale epinglee par empreinte, systeme de fichiers en lecture seule, suppression des capacites Linux. Comparaison des resultats Trivy avant et apres durcissement. Demonstration qu'un controle en echec empeche le deploiement. Procedure de retour arriere demontree.
\end{constat}

% ===========================================================================
\chapter{Supervision, journalisation et reponse aux incidents}
\label{chap:supervision}
% ===========================================================================

\begin{constat}{chapitre a rediger --- phase 7}
P5 supervision continue avec Prometheus et Grafana ; P18 journalisation de securite ; P28 gestion des incidents de securite ; P35 alertes. Evenements de securite journalises conformement a l'ASVS chapitre V16 : issues d'authentification, refus d'autorisation, tentatives de contournement. Expedition des journaux hors de l'hote applicatif. Tableaux de bord et regles d'alerte. Demonstration d'une alerte effectivement declenchee par une attaque simulee. Procedure de reponse aux incidents.
\end{constat}

% ===========================================================================
\chapter{Augmentation par l'intelligence artificielle}
\label{chap:ia}
% ===========================================================================

\begin{constat}{chapitre a rediger --- phase 8}
Deux integrations distinctes. Dans l'application : generateur de description d'annonce fonde sur Gemini, ce qui place le Top 10 OWASP pour les applications LLM dans le perimetre. Dans la pipeline : triage assiste par IA des resultats de securite, publie en commentaire de pull request, sans action autonome. Controles LLM01 injection d'invite, LLM02 traitement non sur de la sortie, LLM06 divulgation d'informations sensibles, LLM10 consommation non bornee. Gouvernance du code genere par IA conformement a la section 7 du sujet. Registre ADR 0009.
\end{constat}

% ===========================================================================
\chapter{Vulnerabilites detectees et corrections}
\label{chap:vulnerabilites}
% ===========================================================================

\begin{constat}{chapitre a rediger --- phase 9}
Chapitre central pour la demonstration exigee par le sujet. Sept vulnerabilites introduites volontairement, chacune sur sa propre branche, chacune detectee par un controle nomme, chacune corrigee. V1 injection SQL detectee par SAST. V2 secret code en dur detecte par Gitleaks, en utilisant le motif de secret JWT et non une cle Google pour les motifs exposes au chapitre 5. V3 dependance vulnerable detectee par analyse de composition. V4 en-tetes de securite supprimes detectes par DAST. V5 controle d'acces rompu detecte par notre propre suite de tests, alors que SAST et DAST ne le detectent pas. V6 image de base obsolete detectee par Trivy. V7 injection d'invite detectee par test dedie. Pour chacune : la chaine complete pratique, outil, integration, detection, resultat, action corrective.
\end{constat}

% ===========================================================================
\chapter{Analyse de l'impact sur la performance de la pipeline}
\label{chap:performance}
% ===========================================================================

\begin{constat}{chapitre a rediger --- phase 10}
Exigence explicite du sujet. Mesure de la duree moyenne sur au moins cinq executions par configuration, depuis la pipeline de reference jusqu'a la pipeline completement outillee, avec calcul du surcout de chaque etage. Identification des etages les plus couteux. Optimisations appliquees et mesurees : mise en cache des dependances, parallelisation des taches, separation entre analyse de reference en pull request et analyse complete nocturne. Comparaison avant et apres optimisation. Compromis retenus et leur justification. Registre ADR 0010.
\end{constat}

% ===========================================================================
\chapter{Limites et perspectives}
\label{chap:limites}
% ===========================================================================

\begin{constat}{chapitre de synthese}
Limites assumees du travail : application mobile reportee en seconde phase avec le detail des travaux restants et de l'analyse du Top 10 OWASP Mobile ; pratiques P19 IAST et P25 RASP non mises en oeuvre, avec justification ; paiement simule et non reel ; environnement de preproduction et non de production. Perspectives : industrialisation, extension du perimetre ASVS au niveau 2 complet, automatisation de la verification ASVS, maturation de la fonction de triage par IA.
\end{constat}

% ===========================================================================
\chapter{Conclusion}
\label{chap:conclusion}
% ===========================================================================

\begin{constat}{chapitre de synthese}
Synthese de la demarche et des resultats. Bilan de la couverture des onze pratiques de niveau Tier 1 et des pratiques de niveau Tier 2 retenues. Retour sur les deux choix structurants du projet : la mise en place de la gouvernance de securite avant tout code applicatif, et la demonstration de l'efficacite des controles par introduction volontaire et controlee de vulnerabilites. Enseignements tires, notamment sur les limites de l'outillage automatise en matiere de controle d'acces et de logique metier.
\end{constat}


% ===========================================================================
%  BIBLIOGRAPHIE
% ===========================================================================
\bibliographystyle{plain}
\bibliography{bibliographie}
\addcontentsline{toc}{chapter}{Bibliographie}

\end{document}
